\chapter{Writing It Up, Publishing and Presenting}\label{ch:writeup}

\weeks{21--24} \quad Milestone M5

Results that nobody can read do not exist. This chapter covers how to write the
project up as a paper-style report, where it could go, how to release the
repository, and how to turn it into CV lines and interview stories. Do not leave
writing to the end: each milestone already produces material for one section.

% =========================================================================
\section{Write as you go}\label{sec:writeup:asyougo}

\begin{table}[htbp]
\centering\small
\caption{Which milestone feeds which section of the report.}
\label{tab:writeup:sections}
\begin{tabularx}{\textwidth}{@{}lL@{}}
\toprule
Milestone & Draft this section right after reaching it \\
\midrule
M0 & Experimental setup: data, metrics, float baseline \\
M1 & Algorithm--hardware co-design: frame study, precision study, the specification \\
M2 & Architecture and verification: block diagram, frame schedule, bit-exactness \\
M3 & Implementation and results on the FPGA: latency, resources, power, demo \\
M4 & Track results: ML accelerator, or ASIC implementation \\
M5 & Introduction, discussion, conclusion, abstract --- written \emph{last} \\
\bottomrule
\end{tabularx}
\end{table}

% =========================================================================
\section{Structure of the report}\label{sec:writeup:structure}

A hardware-for-signal-processing paper usually follows this shape:
\begin{enumerate}
  \item \textbf{Introduction.} The problem (intent-controlled hearing-aid DSP
        must fit a tiny power budget), the gap (the thesis ran on a laptop), the
        contributions in three or four bullet points, and the results in one
        sentence each.
  \item \textbf{Background.} IMCRA, the parametric Wiener filter and the thesis
        system, all brief and cited.
  \item \textbf{Workload characterisation.} The \cref{ch:workload} analysis in
        condensed form. Reviewers love a cost model.
  \item \textbf{Co-design.} Frame configuration, precision and approximations,
        and the gain-LUT control interface, each with the experiment that
        justified it.
  \item \textbf{Architecture and verification.}
  \item \textbf{Results.} FPGA, plus ML or ASIC, against baselines, with
        confidence intervals.
  \item \textbf{Discussion.} What flipped under hardware constraints, threats to
        validity (\cref{ch:evaluation}), and future work, including the
        perceptual calibration study.
  \item \textbf{Conclusion.}
\end{enumerate}

Figures you must have:
\begin{itemize}
  \item the two-plane architecture;
  \item the block diagram;
  \item the frame schedule;
  \item quality against bits;
  \item quality against latency (Pareto);
  \item area or resources against bits;
  \item the headline results table;
  \item a photo of the demo setup.
\end{itemize}

\begin{tipbox}[Writing advice that works]
Write one message per paragraph, and put it in the first sentence. Put the old,
known information at the start of a sentence and the new information at the
end~\cite{ch12-13:gopen1990science}. Write the outline first, then the figures,
then the text~\cite{ch12-13:whitesides2004writing}. For computer-science style
writing in general, see Zobel~\cite{ch12-13:zobel2014writing}. Before you
start, answer the Heilmeier questions: what are you trying to do, how is it done
today, what is new, and who cares~\cite{ch12-13:heilmeier}. They make an
excellent first page.
\end{tipbox}

% =========================================================================
\section{Where it could go}\label{sec:writeup:venues}

A well-executed project of this kind can be submitted to a student-friendly
venue. Check each venue's current call for papers, deadlines and page limits:
\begin{itemize}
  \item \textbf{IEEE NorCAS} (Nordic Circuits and Systems Conference): a
        regional conference with a natural fit for a Danish circuits project.
  \item \textbf{IEEE SiPS} (Workshop on Signal Processing Systems): exactly the
        signal-processing-hardware intersection.
  \item \textbf{IEEE ISCAS}, \textbf{EUSIPCO} or \textbf{ICASSP}: larger, more
        competitive venues. Check whether there are student or late-breaking
        tracks.
  \item \textbf{arXiv}: post a preprint once your supervisor agrees.
\end{itemize}
Within your university, the project is also an excellent basis for an MSc
thesis proposal, or can be formalised as a project course. Ask your study
administration which formats exist for your programme~\cite{ch12-13:dtuprojectcourses}.

% =========================================================================
\section{Release the repository}\label{sec:writeup:release}

\begin{itemize}
  \item A README with the results table at the top, one-command reproduction,
        and the demo video linked.
  \item A tagged release (\code{v1.0}). Archive it for a DOI if you want it to
        be citable, and add a \file{CITATION.cff} file~\cite{ch12-13:githubcitation}.
  \item Licence, data instructions (how to obtain VoiceBank+DEMAND), and pinned
        tool versions.
\end{itemize}

% =========================================================================
\section{CV lines and interview stories}\label{sec:writeup:cv}

Write CV lines with \emph{your measured numbers}. Never fill them in before you
have measured them. Templates:
\begin{itemize}
  \item \enquote{Designed and verified a real-time speech-enhancement datapath
        (IMCRA + parametric Wiener) in SystemVerilog, bit-exact with a Python
        golden model across [N] utterances; [X] ms algorithmic latency, [Y]
        LUTs / [Z] BRAMs on a Zynq-7020.}
  \item \enquote{Reduced per-bin state memory by [k]$\times$ through log-domain
        minimum tracking and frame reconfiguration at a PESQ cost of [d].}
  \item \enquote{Built an int8 matrix--vector accelerator for intent
        classification, preserving macro-F1 within [d] of float.} (Track A)
  \item \enquote{Took a DSP block through an open-source RTL-to-GDS flow
        (LibreLane, [PDK]); [A] mm$^2$, [E] \textmu J/frame from gate-level
        activity.} (Track B)
\end{itemize}

Prepare five stories of about two minutes each. Each one is a situation, what
you did, and what you measured:
\begin{enumerate}
  \item \textbf{A decision that flipped.} mpnet versus MiniLM, or zero padding:
        right for a laptop, wrong for a chip, and how you showed it
        (\cref{ch:workload}).
  \item \textbf{The interface trick.} How language control became a table write
        and two registers.
  \item \textbf{Bit-exact verification.} Why exact equality, and how the golden
        model made debugging systematic.
  \item \textbf{A hard bug.} The first divergence, how you localised it, and the
        regression test you added.
  \item \textbf{A trade-off with numbers.} Latency against quality against
        memory, from your Pareto plot.
\end{enumerate}

% =========================================================================
\section{Talking to industry}\label{sec:writeup:industry}

Hearing-aid manufacturers in Denmark, such as Demant (Oticon), GN and WS
Audiology, build exactly the kind of low-power DSP and ML systems this project
studies. Oticon also has a dedicated research centre,
Eriksholm~\cite{ch12-13:eriksholm}. Your thesis collaboration with Oticon is a
natural first contact: share the demo video and the results table with your
thesis contacts, and ask for twenty minutes of feedback. Engineers respond well
to specific, measured work and a clear question.

\begin{donebox}[Milestone M5 is reached when]
\begin{checklist}
  \item The report is complete, and your supervisor has given feedback that you have addressed.
  \item The repository is released as \code{v1.0}, and reproduces the headline results from one command.
  \item The demo video and results table are linked in the README.
  \item The CV lines contain only measured numbers.
  \item You have rehearsed the five interview stories out loud.
\end{checklist}
\end{donebox}
